This is a small study. We use one width grid (24 values), one sample size $n=200$ and one input dimension per dataset, so the conclusions say nothing about the scaling of peak height with $n$ or about how the effects depend on the aspect ratio $d/n$; the asymptotic theory~\cite{mei2019generalization} is not tested. Only 5 seeds are used, and the heavy tails of the MinNorm peak make the registered paired test on the hump underpowered; the smallest attainable one-sided Wilcoxon $p$ is 0.031. The peak window $0.4\le N/n\le4$, the 5\% hump threshold and the grid for $\lambda$ were chosen by us, the window after one smoke run; a different window would change \texttt{peak\_mse} and \texttt{peak\_pos} for large $\lambda$, where there is no true peak. The grid resolution near the threshold is $0.05$. Our label noise is Gaussian regression noise, not label flips, on both datasets; the digits test labels are clean one-hot vectors and MSE is not accuracy (we do not report accuracy in the paper). The synthetic teacher is outside the model class, which may produce a peak without label noise, and we did not test this. The baselines are our own reimplementations of standard estimators with a closed-form solver (no SGD or early stopping), and the fixed-ridge value $10^{-2}$ was not tuned. The numerical rank threshold for the pseudo-inverse ($10^{-10}s_{\max}$) can matter at $N=n$ where the smallest singular value is tiny; we did not sweep it. Everything ran on CPU in about half an hour; a full study would add several $n$ and $d$, more seeds, real label-flip noise, other activations and non-isotropic data, and a comparison to the asymptotic formulas.
